This section provides the technical documentation for the non-DL physical
priors (\texttt{Try 78} and \texttt{Try 79}) and the final joint architecture
(\texttt{Try 80}). It serves as a bridge between the high-level attempt history
and the numerical results, explaining the exact mathematical transformations,
how every parameter was fitted in practice, and what is genuinely new in the
final \texttt{Try 80} approach.

\begin{table}[H]
\centering
\small
\begin{tabularx}{\textwidth}{p{2.2cm} p{3.1cm} p{4.0cm} X}
\toprule
\textbf{Block} & \textbf{Literature basis} & \textbf{CKM adaptation} &
\textbf{Role in \texttt{Try 80}} \\
\midrule
Path-loss LoS & FSPL and enhanced two-ray propagation &
Effective height-dependent reflection parameters and radial residuals fitted
from CKM LoS pixels &
Frozen LoS prior, injected as an input channel and used as the anchor for
bounded residual learning. \\
\addlinespace
Path-loss NLoS & COST-231/Hata-style urban loss and A2G elevation envelopes &
Regime-wise calibration with local density, height, shadowing and elevation
features &
Frozen NLoS prior selected by the explicit LoS/NLoS geometry mask. \\
\addlinespace
Delay and angular spread & Log-domain large-scale spread modelling and
LoS/NLoS separation &
Regime-wise ridge regression in \(\log(1+y)\) space using multiscale
morphology and obstruction features &
Frozen spread priors used as anchors for the multi-task residual model. \\
\addlinespace
\texttt{Try 80} residual model & Probabilistic regression and mixture-density
ideas &
Prior-anchored per-pixel GMM with residual caps, anchor gates, explicit
region masks and shared multi-task trunk &
Learns residual uncertainty and local corrections without replacing the
interpretable priors. \\
\bottomrule
\end{tabularx}
\caption{One-page summary of what is inherited from the literature, what is
adapted to CKM, and what is reused or newly assembled in \texttt{Try 80}.}
\label{tab:executive_summary}
\end{table}

\subsection{High-level architecture}

The central design decision in \texttt{Try 80} is \emph{not} to learn path
loss, delay spread, and angular spread from scratch. Instead:

\begin{itemize}
\item \texttt{Try 78} provides a frozen prior for \texttt{path\_loss},
\item \texttt{Try 79} provides frozen priors for \texttt{delay\_spread} and
      \texttt{angular\_spread},
\item \texttt{Try 80} learns bounded residual corrections around those priors.
\end{itemize}

The deterministic or interpretable structure is kept in the prior, and the deep
model is used only for the residual part that the prior cannot explain.
Figure~\ref{fig:pipeline} illustrates the overall pipeline.

\begin{figure}[H]
\centering
\begin{tikzpicture}[
  box/.style={rectangle, draw, rounded corners=3pt, minimum width=2.2cm,
              minimum height=0.75cm, align=center, font=\footnotesize},
  arr/.style={-{Stealth[length=5pt]}, thick},
  node distance=0.5cm and 1.1cm
]
\node[box, fill=blue!15]  (t78) {Try 78\\(path loss)};
\node[box, fill=green!15, below=of t78] (t79) {Try 79\\(spreads)};
\node[box, fill=orange!15, right=1.2cm of t78, yshift=-0.62cm] (inject)
      {Prior injection\\(input channels)};
\node[box, fill=red!15, right=1.2cm of inject] (dnn)
      {DNN residual\\multi-task};
\node[box, fill=gray!20, right=1.2cm of dnn] (out)
      {PL, DS, AS\\outputs};

\draw[arr] (t78) -- (inject);
\draw[arr] (t79) -- (inject);
\draw[arr] (inject) -- (dnn);
\draw[arr] (dnn) -- (out);

\node[below=0.25cm of dnn, font=\scriptsize, text width=2.4cm, align=center]
      {$\delta \in [-\delta_{\max}, +\delta_{\max}]$};
\end{tikzpicture}
\caption{Overall \texttt{Try 80} pipeline. The frozen priors are injected into the DNN, which learns a bounded correction $\delta$ around them.}
\label{fig:pipeline}
\end{figure}

\subsection{Prior design rules}
\label{sec:good_prior_design}

The main lesson from the prior phase is that a useful prior is not just a
physical-looking equation. It is a complete estimator: the formula, fitted
calibration, feature units, LoS/NLoS mask, topology features, and evaluation
mask must be identical between fitting and inference. Earlier prior attempts
used many of the same physical ingredients as \texttt{Try 78}, but they did
not always keep those pieces aligned. \texttt{Try 78} became strong because it
made the prior self-contained and auditable under the same city-holdout
contract used for the final model.

The distinction matters most for NLoS. A calibration fitted for one physical
formula cannot reliably correct another; local height and density features
must have the same units when fitted and when applied; and regime assignment
should come from the sample geometry rather than only from a broad city label.
Once those constraints were enforced, the prior could separate deterministic
large-scale structure from the genuinely local residual that Try~80 later
learned.

This does not mean that all earlier NLoS failures were only implementation
mismatches. The no-prior and weak-prior neural branches also failed because
single-output point regression was a poor match for blocked, multi-modal,
heavy-tailed NLoS and spread targets. That is why the thesis keeps both
lessons: \texttt{Try 78}/\texttt{Try 79} fix the prior-estimator problem, while
\texttt{Try 76}/\texttt{Try 77} motivate distribution-aware heads and residual
uncertainty.

\begin{table}[H]
\centering
\caption{Compact design rules extracted from the prior attempts.}
\label{tab:prior_design_rules}
\small
\begin{tabularx}{\textwidth}{>{\raggedright\arraybackslash}p{0.28\textwidth}
>{\raggedright\arraybackslash}X}
\toprule
Rule & Effect on the final pipeline \\
\midrule
Fit and apply the same formula &
The calibration corrects the actual runtime prior, not a neighbouring variant. \\
Keep units explicit &
Distances, heights, topology values, and normalized features have one shared
definition across fitting, validation, and test. \\
Use the metric mask during fitting &
The prior is optimized for valid receiver pixels rather than for building cells
that are excluded from the final metric. \\
Hard-separate LoS and NLoS &
The two regimes keep different physics, feature sets, and calibration
parameters when the dataset provides a reliable mask. \\
Let literature define axes, not all coefficients &
Distance, height, frequency, and elevation come from propagation theory;
effective reflection and local morphology corrections are fitted from CKM. \\
Learn bounded residuals after the prior &
Try~80 is allowed to correct local errors without overwriting the calibrated
large-scale structure. \\
\bottomrule
\end{tabularx}
\end{table}

\subsection{Literature ingredients and thesis adaptations}

The final pipeline is a hybrid between classical propagation models,
distribution-aware regression, and dataset-specific calibration. Table
~\ref{tab:prior_literature_mapping} keeps the mapping explicit; the detailed
equations and fitting procedures follow in the dedicated \texttt{Try 78},
\texttt{Try 79}, and \texttt{Try 80} subsections.

\begin{table}[H]
\centering
\caption{Literature-to-design mapping used by the final prior-anchored model.}
\label{tab:prior_literature_mapping}
\scriptsize
\renewcommand{\arraystretch}{1.12}
\begin{tabularx}{\textwidth}{>{\raggedright\arraybackslash}p{0.18\textwidth}
>{\raggedright\arraybackslash}X
>{\raggedright\arraybackslash}X
>{\raggedright\arraybackslash}X}
\toprule
Block & Literature basis & CKM adaptation & New role in this thesis \\
\midrule
Try~78 LoS path loss &
FSPL and coherent two-ray propagation. &
Effective complex reflection, height-bin interpolation, and radial residual
fitting on CKM LoS pixels. &
Reconstructs the radial interference pattern and supplies the frozen LoS
anchor. \\
Try~78 NLoS path loss &
COST-231/Hata-style attenuation, A2G elevation envelopes, and LoS/NLoS
separation. &
Map-derived density, height, NLoS-support, elevation, and shadowing features
with regime-wise calibration. &
Turns a broad empirical curve into a local, topology-aware NLoS prior. \\
Try~79 delay and angular spread &
Log-normal spread statistics from 3GPP/WINNER-style channel models and UAV
elevation trends. &
Log-domain ridge regression with local morphology features and a sparse-regime
fallback hierarchy. &
Provides stable spread anchors without pretending that spreads are smooth
deterministic path-loss maps. \\
Try~80 residual model &
U-Net-style dense prediction, FiLM conditioning, probabilistic regression, and
mixture-density ideas. &
Prior channels, explicit masks, residual caps, anchor gates, and a shared
multi-task trunk. &
Learns only the bounded local correction and uncertainty left after the
calibrated priors. \\
\bottomrule
\end{tabularx}
\end{table}
